\chapter{Hỏi Để Kết Thúc Một Ngày}
\chapterintro{Một câu cuối ngày có thể giữ lại cả tinh thần lớp}{Kết một ngày học gọn và đúng giúp ngày mai vào nhịp dễ hơn.}

\section{Hôm nay em muốn giữ lại điều gì nhất?}
\begin{cauhoibox}{Câu hỏi}
Hôm nay em muốn giữ lại điều gì nhất?
\end{cauhoibox}
\begin{taisaocoich}
Câu hỏi này giúp học sinh nhặt lại một điểm sáng trong ngày, dù ngày đó bận hay mệt.
\end{taisaocoich}
\begin{goiybien}
Hợp cuối buổi, cuối ngày, hoặc cuối tuần. Có thể trả lời bằng một từ, một câu ngắn, hoặc giơ tay chọn.
\end{goiybien}
\ngancach

\section{Ngày hôm nay làm em mệt ở chỗ nào và mạnh hơn ở chỗ nào?}
\begin{cauhoibox}{Câu hỏi}
Ngày hôm nay làm em mệt ở chỗ nào và mạnh hơn ở chỗ nào?
\end{cauhoibox}
\begin{taisaocoich}
Đây là câu hỏi hai mặt: vừa nhìn nhận khó khăn, vừa không bỏ quên điều tích cực học sinh đã có.
\end{taisaocoich}
\begin{goiybien}
Có thể cho lớp viết nhanh vào giấy nếu không muốn phát biểu công khai.
\end{goiybien}
\ngancach

\section{Nếu ngày mai bắt đầu lại, em muốn bắt đầu khác ở điểm nào?}
\begin{cauhoibox}{Câu hỏi}
Nếu ngày mai bắt đầu lại, em muốn bắt đầu khác ở điểm nào?
\end{cauhoibox}
\begin{taisaocoich}
Câu hỏi này hướng học sinh sang chuyển động tiếp theo, không mắc kẹt trong cảm giác của hôm nay.
\end{taisaocoich}
\begin{goiybien}
Đây là câu rất đẹp để chốt cuối tuần hoặc cuối ngày nhiều áp lực.
\end{goiybien}
\ngancach

\section{Ai hoặc điều gì đáng cảm ơn trong hôm nay?}
\begin{cauhoibox}{Câu hỏi}
Ai hoặc điều gì đáng cảm ơn trong hôm nay?
\end{cauhoibox}
\begin{taisaocoich}
Cuối ngày, một câu hỏi cảm ơn nhỏ giúp học sinh nhặt lại điểm sáng thay vì chỉ mang sự mệt về nhà. Nó làm nhịp kết mềm hơn rất nhiều.
\end{taisaocoich}
\begin{goiybien}
Cho trả lời bằng một từ, một tên người, hoặc một khoảnh khắc. Không cần phân tích dài.
\end{goiybien}
\ngancach

\section{Nếu gửi cho ngày mai một lời nhắn ngắn, em sẽ nhắn gì?}
\begin{cauhoibox}{Câu hỏi}
Nếu gửi cho ngày mai một lời nhắn ngắn, em sẽ nhắn gì?
\end{cauhoibox}
\begin{taisaocoich}
Câu hỏi này giúp học sinh kết ngày hôm nay bằng một hướng đi tiếp. Nó vừa mang tính tự nhìn lại, vừa tạo một chuyển động nhỏ cho ngày mới.
\end{taisaocoich}
\begin{goiybien}
Rất hợp cuối ngày hoặc cuối tuần. Có thể cho lớp nói nối tiếp thật nhanh để tạo một đoạn kết đẹp.
\end{goiybien}
